% eksperymentalne tłumaczenie na włoski

%

\chapter{Współliniowość} % lang-pl
\chapter{Allineamento} % lang-it


Opiszemy najpierw trzy twierdzenia (Menelausa, Cevy, van Aubela), a potem cztery inne twierdzenia (Desargues'a, Pappusa, Pascala, Brianchona). % lang-pl
Descriveremo dapprima tre teoremi (Menelao, Ceva, van Aubel), quindi altri quattro teoremi (Desargues, Pappo, Pascal, Brianchon). % lang-it


Twierdzenie Menelaosa zostanie odkryte jeszcze w starożytności, zaś Cevy dopiero w XVII wieku, mimo że obydwa będą zazwyczaj wykładane razem w późniejszych podręcznikach geometrii, np. u Evesa \cite[s. 63-67]{eves1_1972}, Zetela \cite[s. 13, 43]{zetel_2020}, % lang-pl
Il teorema di Menelao fu scoperto già nell'antichità, mentre quello di Ceva soltanto nel XVII secolo, sebbene entrambi vengano solitamente presentati insieme nei successivi manuali di geometria, ad esempio da Eves \cite[pp. 63-67]{eves1_1972} e Zetel \cite[pp. 13, 43]{zetel_2020}. % lang-it

Mają liczne zastosowania: % lang-pl
Essi hanno numerose applicazioni: % lang-it

prostsze (wysokości, albo środkowe, albo dwusieczne kątów w trójkącie przecinają się w jednym punkcie) albo trudniejsze: % lang-pl
più semplici (le altezze, oppure le mediane, oppure le bisettrici di un triangolo si incontrano in un unico punto) oppure più complesse: % lang-it

w dowodzie twierdzenia Desarguesa o dwóch trójkątach (dwa trójkąty są współbiegunowe wtedy i tylko wtedy, gdy są współosiowe), % lang-pl
nella dimostrazione del teorema di Desargues sui due triangoli (due triangoli sono prospettici da un punto se e solo se sono prospettici da una retta), % lang-it

albo w dowodzie twierdzenia Pascala o heksagramie (punkty przecięcia przeciwległych krawędzi sześciokąta wypukłego albo nie, ale wpisanego w okrąg, leżą na jednej prostej\footnote{Sześć punktów na okręgu tworzy 60 sześciokątów, które wyznaczają 60 prostych Pascala, przechodzących po trzy przez 20 punktów Steinera, które leżą na 15 prostych Plueckera. Są też punkty Kirkmana, proste Cayleya i punkty Salmona (!)}). % lang-pl
oppure nella dimostrazione del teorema dell'esagramma di Pascal (i punti d'intersezione dei lati opposti di un esagono, convesso oppure no, ma inscritto in un cerchio, giacciono su una stessa retta\footnote{Sei punti su un cerchio determinano 60 esagoni, che individuano 60 rette di Pascal; queste passano a gruppi di tre per 20 punti di Steiner, che giacciono su 15 rette di Plücker. Vi sono inoltre i punti di Kirkman, le rette di Cayley e i punti di Salmon (!)}). % lang-it

% TODO: rewrite
Wprowadzenie do geometrii punktów w nieskończoności zwykle przypisuje się Johannowi Keplerowi (1571..1630), jednak to Gerard Desargues (1593..1662) w opracowaniu dotyczącym krzywych stożkowych (jego Brouillon projet), opublikowanym w 1639 roku, jako pierwszy zastosował tę ideę w sposób systematyczny. Praca Desargues’a stanowi pierwszy istotny postęp w geometrii syntetycznej od czasów starożytnych Greków. % lang-pl

L'introduzione dei punti all'infinito in geometria viene solitamente attribuita a Johannes Kepler (1571--1630), ma fu Gérard Desargues (1593--1662), nel suo trattato sulle coniche (\emph{Brouillon projet}), pubblicato nel 1639, a utilizzare per primo questa idea in modo sistematico. L'opera di Desargues rappresenta il primo progresso significativo della geometria sintetica dai tempi degli antichi Greci. % lang-it

% eksperymentalne tłumaczenie na włoski

%

\section{Twierdzenie Menelausa} % lang-pl
\section{Teorema di Menelao} % lang-it

Znamy trzy twierdzenia o współliniowości: o prostej Eulera, o prostej Wallace'a-Simsona, o prostej Auberta. % lang-pl
Conosciamo tre teoremi di allineamento: sulla retta di Eulero, sulla retta di Wallace-Simson e sulla retta di Aubert. % lang-it

\begin{proposition}
[twierdzenie Menelaosa] % lang-pl
[teorema di Menelao] % lang-it
	...
	Wówczas punkty $K, L, M$ są współliniowe wtedy i tylko wtedy, gdy zachodzi % lang-pl
	Allora i punti $K$, $L$ e $M$ sono allineati se e solo se vale % lang-it
	\begin{equation}
		[AMB] [BKC] [CLA] = -1.
	\end{equation}
	\index{twierdzenie!Menelaosa} % lang-pl
	\index{teorema!di Menelao} % lang-it
\end{proposition}

% Eves s. 248

\index[persons]{Menelaos z Aleksandrii}% % lang-pl
\index[persons]{Menelao di Alessandria}% % lang-it

Brak pewności co do tego, kto będzie pierwszym odkrywcą tego. % lang-pl
Ale najstarsze ocalałe do czasów późniejszych wyjaśnienie pojawia się u Menelaosa (Σφαιρική) jako lemat, użyteczny w sferycznej wersji twierdzenia. % lang-pl
Non vi è certezza su chi ne sia stato il primo scopritore. % lang-it
Tuttavia, la più antica esposizione giunta fino a noi compare nella \emph{Sphaerica} (Σφαιρική) di Menelao come lemma utile per la versione sferica del teorema. % lang-it

Piszą o nim Guzicki \cite[s. 84]{guzicki_2021}; Hartshorne \cite[s. 180, 181]{hartshorne_2000}; Audin \cite[s. 38]{audin_2003}; Zetel \cite[s. 43]{zetel_2020}. % lang-pl
Uogólnieniem twierdzenia Menelaosa będzie twierdzenie Carnota (udowodnione w 1803 roku w \emph{,,Géométrie de position}''), że iloczyn długości odcinków wyznaczonych przez prostą na bokach wielokąta i niemających wspólnych końców jest równy iloczynowi długości pozostałych odcinków; % lang-pl
Ne scrivono Guzicki \cite[p. 84]{guzicki_2021}; Hartshorne \cite[pp. 180, 181]{hartshorne_2000}; Audin \cite[p. 38]{audin_2003} e Zetel \cite[p. 43]{zetel_2020}. % lang-it
Una generalizzazione del teorema di Menelao è il teorema di Carnot (dimostrato nel 1803 nella \emph{``Géométrie de position''}), secondo cui il prodotto delle lunghezze dei segmenti determinati da una retta sui lati di un poligono e privi di estremi comuni è uguale al prodotto delle lunghezze dei segmenti rimanenti; % lang-it
Pisze o nim Zetel w formie ćwiczenia \cite[s. 45]{zetel_2020}. % lang-pl
Ne scrive Zetel come esercizio \cite[p. 45]{zetel_2020}. % lang-it
% https://encyclopediaofmath.org/wiki/Carnot_theorem
% Zetel 45.

% TODO: https://en.wikipedia.org/wiki/Menelaus%27s_theorem
% It is uncertain who actually discovered the theorem; however, the oldest extant exposition appears in Spherics by Menelaus. In this book, the plane version of the theorem is used as a lemma to prove a spherical version of the theorem.

% TODO: https://www.deltami.edu.pl/2011/03/twierdzenie-menelaosa/

%
%

\section{Twierdzenie Cevy}

Ważnym kryterium współpękowości trzech czewian jest:

\begin{theorem}[Cevy]
\label{twierdzenie_cevy}%
	Dany jest trójkąt $ABC$ i trzy różne od wierzchołków punkty $K \in BC$, $L \in CA$, $M \in AB$.
	Wówczas czewiany $AK$, $BL$, $CM$ są współpękowe wtedy i tylko wtedy, gdy
	\begin{equation}
		[AMB] [BKC] [CLA] = 1.
	\end{equation}
	\index{twierdzenie!Cevy}
\end{theorem}
% Neugebauer 119
% Ceva = guzicki-3
% Eves s. 248

https://www.deltami.edu.pl/2021/09/zabawy-na-polu/ 7

Twierdzenie należy do geometrii afinicznej, ponieważ można je wyrazić bez odnoszenia się do kątów, pól albo długości.
Przypisze się je Giovanniemu Cevie, który wyda pracę \emph{De lineis rectis} (o~liniach prostych) w 1678 roku; chociaż twierdzenie nie będzie obce Yusufowi al-Mu'tamanowi ibn Hudzie, emirowi taify w Saragossie z XI wieku.
\index[persons]{al-Mutaman, Yusuf ibn Huda}%
\index[persons]{Ceva, Giovanni}%
Wspomniana praca Cevy będzie zawierać także twierdzenie Menelaosa, odkryte na nowo przez Cevę.

Piszą o nim (twierdzeniu, nie emirze) Zetel \cite[s. 11-14, 40]{zetel_2020}, Audin \cite[s. 38]{audin_2003}, Hartshorne \cite[s. 182]{hartshorne2000}.
Zetel stwierdzi, że uogólnienie twierdzenia Cevy nazywa się twierdzeniem Ponceleta, mimo braku podobieństw do faktu \ref{big_poncelet}.
Osiem zadań, do rozwiązania których wygodnym narzędziem jest twierdzenie Cevy, poda Joanna Jaszuńska w $\Delta_{11}^2$. % https://www.deltami.edu.pl/2011/02/twierdzenie-cevy/

Wnioskiem z twierdzenia Cevy jest wynik znany jako lemat Steinbarta:

\begin{lemma}[Steinbarta]
	W trójkąt $ABC$ wpisano okrąg, styczny do boków $BC$, $AC$ i $AB$ kolejno w punktach $D$, $E$, $F$.
	Punkty $X$, $Y$, $Z$ leżą odpowiednio na krótszych łukach $EF$, $FD$ i $DE$.
	Następujące warunki są równoważne:
	\begin{itemize}
		\item proste $XD$, $YE$, $ZF$ są współpękowe,
		\item proste $AX$, $BY$, $CZ$ są współpękowe.
	\end{itemize}
\end{lemma}

Piszą o nim Ambroszczyk, Komisarczyk \cite[s. 17]{komisarczyk_2024}, ale pełną historię wyniku znajdziemy w artykule
% https://www.tandfonline.com/doi/full/10.1080/0025570X.2025.2595954#d1e168

% @article{Humenberger01012026,
% author = {Hans Humenberger},
% title = {An Introduction to and a Further Variation of the Steinbart Theorem},
% journal = {Mathematics Magazine},
% volume = {99},
% number = {1},
% pages = {20--29},
% year = {2026},
% publisher = {Taylor \& Francis},
% doi = {10.1080/0025570X.2025.2595954},
% URL = { 
%         https://doi.org/10.1080/0025570X.2025.2595954
% },
% eprint = { 
%         https://doi.org/10.1080/0025570X.2025.2595954
% }
% }

% * The phenomenon described in the Steinbart theorem was already formulated 1991 as a problem by S. Rabinowitz ([4], of course, not calling it Steinbart theorem; only for points P in the interior of the triangle) and a solution was given by F. Bellot and M. A. Lopez, in the following year 1992, see [5]. If a reader happens to know an even earlier source of the phenomenon described in the Steinbart theorem, please inform the author.
% https://www.tandfonline.com/doi/full/10.1080/0025570X.2025.2595954#abstract
% The Steinbart theorem seems to be not very well known in the community, it was detected and proved by an ambitious German high school student, Oliver Funck.
% The high school student Oliver Funck (Steinbart gymnasium in Duisburg, Germany) detectedFootnote* and proved (with the help of a computer algebra system; below we present a geometric proof) this phenomenon in the year 2000 presenting a note (16 pp) in the framework of “Jugend forscht” (German for “youth researches,” see [Citation7]). To the honor of the school’s eponym, Quintin Steinbart, he called Q the Steinbart point of the triangle ABC with respect to point P and the corresponding theorem Steinbart’s theorem. This labeling should remind one of the Exeter point, a special case of the Steinbart point, namely 𝑃= centroid G of the triangle, detected in 1986 at a computers-in-mathematics workshop at Phillips Exeter Academy (this point lies on the Euler line). Since the Exeter point was named after the institution in which it was detected, Oliver Funck also named “his” point after “his” institution (Steinbart gymnasium).


Prostą, która przechodzi przez wierzchołek trójkąta i przeciwległy bok nazywa się czewianą, nazwy tej używa się też wobec odcinka trójkąta leżącego na tej prostej.

\begin{proposition}[wzór Steinera-Routha]
	Rozpatrzmy trójkąt $ABC$ z czewianami $AD$, $BE$, $CF$ tak, że
	\begin{equation}
		\frac{BD}{DC} = \lambda_1, \quad
		\frac{CE}{EA} = \lambda_2, \quad
		\frac{AF}{FB} = \lambda_3.
	\end{equation}
	Oznaczmy przekrój $BE$ i $CF$ przez $G_1$; $AD$ i $CF$ przez $G_2$, zaś $AD$ i $BE$ przez $G_3$.
	Wtedy
	\begin{equation}
		\frac{S_{G_1G_2G_3}}{S_{ABC}} = \frac{(\lambda_1 \lambda_2 \lambda_3 - 1)^2}{
			(\lambda_1 \lambda_2 + \lambda_1 + 1)
			(\lambda_2 \lambda_3 + \lambda_2 + 1)
			(\lambda_3 \lambda_1 + \lambda_3 + 1)
		}
	\end{equation}
	oraz
	\begin{equation}
		\frac{S_{DEF}}{S_{ABC}} = \frac{\lambda_1 \lambda_2 \lambda_3 + 1}{
			(\lambda_1 + 1)
			(\lambda_2 + 1)
			(\lambda_3 + 1)
		}.
	\end{equation}
\end{proposition}

% J. Steiner, Gesammelte Werke, Vol. 1, Reimer, Berlin, 1882. strony 163-168
% E. J. Routh, A Treatise on Analytical Statics with Numerous Examples, Vol. 1, second edition. Cambridge University Press, 1896. strona 82
% https://ems.press/content/serial-article-files/51635

Stanowi to uogólnienie twierdzeń Cevy ($\lambda_1 \lambda_2 \lambda_3 = 1$) i Menelausa ($\lambda_1 \lambda_2 \lambda_3 = -1$).
Napiszą o nim w Delcie ($\Delta_{11}^6$, $\Delta_{12}^3$, $\Delta_{16}^8$)

\begin{example}[trójkąt, który zaskoczył Feynmana]
	Dla $\lambda_1 = \lambda_2 = \lambda_3 = 2$, stosunek pól trójkątów wynosi $1 : 7$.
\end{example}

% https://www.deltami.edu.pl/2011/06/sprawa-niezbyt-pedagogiczna/
% https://www.deltami.edu.pl/2012/03/co-moga-nam-dac-ciezary-i-wypory/
% https://www.deltami.edu.pl/2016/08/polecamy-bardzo-trudne-zadanie/

Ceva 3D https://www.deltami.edu.pl/2013/10/kacik-przestrzenny-19-plaszczyzny-przecinajace-sie-w-punkcie/

\subsubsection{Punkt Nagela}
% eksperymentalne tłumaczenie na włoski

%

TODO: Punkt Nagela, pojawia się u Evesa \cite[s. 71]{eves1_1972} jako ćwiczenie na zastosowanie twierdzenia Cevy razem z punktem Gergonne'a. % lang-pl

\index{punkt!Nagela} % lang-pl
TODO: Il punto di Nagel compare in Eves \cite[p. 71]{eves1_1972} come esercizio sull'applicazione del teorema di Ceva insieme al punto di Gergonne. % lang-it

\index{czewiany Nagela} % lang-pl
\index{punto!di Nagel} % lang-it

\index{ceviane di Nagel} % lang-it

\begin{proposition}[punkt Nagela] % lang-pl
\begin{proposition}[punto di Nagel] % lang-it
    \label{punkt_nagela}
\end{proposition}

O punkcie Nagela pisze też Zetel \cite[s. 22, 25, 64]{zetel_2020} (punkt Nagela, punkt przecięcia środkowych -- środek ciężkości trójkąta, środek okręgu wpisanego oraz środek ciężkości obwodu trójkąta leżą na jednej prostej zwanej drugą prostą Eulera albo prostą Eulera-Nagela). % lang-pl

Del punto di Nagel scrive anche Zetel \cite[p. 22, 25, 64]{zetel_2020} (punto di Nagel, punto di concorrenza delle mediane: il baricentro del triangolo, l'incentro e il baricentro del perimetro del triangolo giacciono su una stessa retta, detta seconda retta di Eulero oppure retta di Eulero--Nagel). % lang-it

% % Punkty izogonalnie sprzężone w trójkącie % https://www.deltami.edu.pl/2010/11/izogonalnie-sprzezone/

\subsubsection{Punkt Gergonne'a}
% eksperymentalne tłumaczenie na włoski

%

TODO: Czewiany Gergonne'a. % lang-pl
\index{czewiany Gergonne'a} % lang-pl
TODO: Ceviane di Gergonne. % lang-it
\index{ceviane di Gergonne} % lang-it

\begin{definition}
[punkt Gergonne'a] % lang-pl
[punto di Gergonne] % lang-it
\index{punkt!Gergonne'a}% % lang-pl
\index{punto!di Gergonne}% % lang-it
\label{punkt_gergonne}
    Dany jest okrąg wpisany w~trójkąt $ABC$, styczny do boków $BC$, $AC$, $AB$ odpowiednio w~punktach $P$, $Q$ i $R$. % lang-pl
    Wtedy czewiany $AP$, $BQ$ i $CR$ są współpękowe. % lang-pl
    Sia dato il cerchio inscritto nel triangolo $ABC$, tangente ai lati $BC$, $AC$, $AB$ rispettivamente nei punti $P$, $Q$ e $R$. % lang-it
    Allora le ceviane $AP$, $BQ$ e $CR$ sono concorrenti. % lang-it
\end{definition}

TODO: rysunek % lang-pl
TODO: figura % lang-it

To samo, ale bez użycia słowa ,,czewiany'' napisze Zetel \cite[s. 15, 25]{zetel_2020}. % lang-pl
Lo stesso risultato, ma senza usare la parola ,,ceviane'', è presentato da Zetel \cite[p. 15, 25]{zetel_2020}. % lang-it

%

\subsubsection{Punkt Lemoine'a}
% eksperymentalne tłumaczenie na włoski

%

\index{punkt!Lemoine'a} % lang-pl
\index{punto!di Lemoine} % lang-it

Emile Lemoine odczyta w 1873 roku pracę \emph{,,Sur quelques propriétés d'un point remarquable du triangle''}, gdzie po raz pierwszy pojawią się bez dowodów fakty \ref{random_prp_symmedians_proportional_squares} i \ref{punkt_lemoine}. % lang-pl
Termin symediana zaproponuje Maurice d'Ocagne w 1883 roku, jako zastępstwo \emph{antyrównoległej środkowej} używanego przez Lemoine'a. % lang-pl
Nel 1873 Émile Lemoine presenterà il lavoro \emph{,,Sur quelques propriétés d'un point remarquable du triangle''}, nel quale compariranno per la prima volta, senza dimostrazione, i fatti \ref{random_prp_symmedians_proportional_squares} e \ref{punkt_lemoine}. % lang-it
Il termine \emph{symmediana} sarà proposto da Maurice d'Ocagne nel 1883 come sostituto di \emph{mediana antiparallela}, espressione usata da Lemoine. % lang-it
% Nouvelles Annales de Mathématiques, 3rd series, II, 451 (1883).
Punkt Lemoine'a pojawia się u Józefa Neuberga (praca \emph{Mémoire sur le tétraèdre} z 1886 roku). % lang-pl
Il punto di Lemoine compare nell'opera di Joseph Neuberg (\emph{Mémoire sur le tétraèdre}, 1886). % lang-it
% https://www.persee.fr/doc/marb_0770-8459_1886_num_37_1_2317
W~1847 roku Ernst Grebe skonstruuje kwadraty na zewnątrz trójkąta i odkryje punkt Lemoine'a przed Lemoinem, dlatego w literaturze padać będzie też nazwa punkt Grebego. % lang-pl
Nel 1847 Ernst Grebe costruirà quadrati all'esterno di un triangolo e scoprirà il punto di Lemoine prima dello stesso Lemoine; per questo motivo nella letteratura compare anche il nome di punto di Grebe. % lang-it

Wszystko to pięknie opisze Altshiller-Court \cite{altshiller_2007}. % lang-pl
Tutto ciò è descritto magnificamente da Altshiller-Court \cite{altshiller_2007}. % lang-it

% UW: Twierdzenie Cevy (wraz z trygonometryczną wersją),
% przykłady punktów szczególnych trójkąta:
% punkt Nagela (Guzicki-4),
% punkt Gergonne'a (guzicki4),

% https://www.deltami.edu.pl/2016/05/izogonalne-sprzezenie-i-symediany/

\begin{definition}
[symediana] % lang-pl
[symmediana] % lang-it
    Odbicie symetryczne środkowej trójkąta względem dwusiecznej wychodzącej z~tego samego wierzchołka nazywamy symedianą. % lang-pl
    La riflessione della mediana di un triangolo rispetto alla bisettrice uscente dallo stesso vertice si chiama symmediana. % lang-it
\end{definition}

Symediana to prosta izogonalna ze środkową (Zetel \cite[s. 111]{zetel_2020}; Ambroszczyk i Komisarczyk \cite[s. 26]{komisarczyk_2024}; Dutka w $\Delta_{15}^2$). % lang-pl
La symmediana è la retta izogonale della mediana (Zetel \cite[p. 111]{zetel_2020}; Ambroszczyk e Komisarczyk \cite[p. 26]{komisarczyk_2024}; Dutka in $\Delta_{15}^2$). % lang-it
% https://www.deltami.edu.pl/2015/02/krotka-opowiesc-o-symedianie/
Symediana jest miejscem geometrycznym punktów, których odległości od dwóch boków trójkąta są proporcjonalne do tych boków. % lang-pl
Ale dużo ważniejszą jej własnością jest: % lang-pl
La symmediana è il luogo geometrico dei punti le cui distanze da due lati del triangolo sono proporzionali a tali lati. % lang-it
Ma una proprietà molto più importante è la seguente: % lang-it

\begin{proposition}
    \label{random_prp_symmedians_proportional_squares}
    Symediana poprowadzona z wierzchołka $C$ trójkąta $ABC$ dzieli wewnętrznie przeciwległy bok proporcjonalnie do kwadratów długości boków przyległych: % lang-pl
    La symmediana condotta dal vertice $C$ del triangolo $ABC$ divide internamente il lato opposto in proporzione ai quadrati delle lunghezze dei lati adiacenti: % lang-it
    \begin{equation}
        \frac{AF}{BF} = \frac{b^2}{a^2}.
    \end{equation}
\end{proposition}

I odwrotnie, punkt leżący na boku $AB$ i spełniający powyższą proporcję wyznacza symedianę z~wierzchołka $C$. % lang-pl
Zetel \cite[s. 111, 112]{zetel_2020} uważa to za szczególny przypadek twierdzenie Steinera o~prostych izogonalnych. % lang-pl
(Chyba nie ma tego twierdzenia w tej książce). % lang-pl
Viceversa, un punto appartenente al lato $AB$ che soddisfa la proporzione precedente determina la symmediana uscente dal vertice $C$. % lang-it
Zetel \cite[p. 111, 112]{zetel_2020} considera questo un caso particolare del teorema di Steiner sulle rette izogonali. % lang-it
(Probabilmente questo teorema non compare nel presente libro). % lang-it

\begin{proposition}
    Symediana wychodząca z wierzchołka $A$ trójkąta $ABC$ ma długość % lang-pl
    La symmediana uscente dal vertice $A$ del triangolo $ABC$ ha lunghezza % lang-it
    \begin{equation}
        \frac{bc}{b^2 + c^2} \cdot \sqrt {2b^2 + 2c^2 - a^2}.
    \end{equation}
\end{proposition}

Zetel \cite[s. 119]{zetel_2020} pisze, że można dostać tę długość z twierdzenia Stewarta. % lang-pl
Zetel \cite[p. 119]{zetel_2020} scrive che questa lunghezza può essere ottenuta mediante il teorema di Stewart. % lang-it
% ale nie trzeba.

\begin{proposition}
[punkt Lemoine'a] % lang-pl
[punto di Lemoine] % lang-it
    \label{punkt_lemoine}%
    W ustalonym trójkącie $ABC$, trzy symediany przecinają się w jednym punkcie $L$, zwanym punktem Lemoine'a. % lang-pl
    In un triangolo fissato $ABC$, le tre symmediane si incontrano in un unico punto $L$, detto punto di Lemoine. % lang-it
\end{proposition}

Zetel \cite[s. 114]{zetel_2020} sugeruje udowodnić \ref{punkt_lemoine} na podstawie twierdzenia odwrotnego do twierdzenia Cevy. % lang-pl
Ross Honsberger nazwie istnienie punktu Lemoine'a ,,jednym z klejnotów koronnych współczesnej geometrii''. % lang-pl
Zetel \cite[p. 114]{zetel_2020} suggerisce di dimostrare \ref{punkt_lemoine} mediante il teorema inverso di Ceva. % lang-it
Ross Honsberger definirà l'esistenza del punto di Lemoine come \emph{uno dei gioielli della corona della geometria moderna}. % lang-it

\begin{proposition}
    Punkt Lemoine'a minimalizuje sumę kwadratów odległości od boków trójkąta. % lang-pl
    Il punto di Lemoine minimizza la somma dei quadrati delle distanze dai lati del triangolo. % lang-it
\end{proposition}

% TODO: https://www.deltami.edu.pl/media/articles/2015/02/delta-2015-02-krotka-opowiesc-o-symedianie.pdf

\begin{proposition}
    Punkt Lemoine'a $L$ trójkąta $ABC$ jest środkiem ciężkości trójkąta, którego wierzchołki to rzuty punktu $L$ na boki $AB$, $BC$, $AC$. % lang-pl
    Il punto di Lemoine $L$ del triangolo $ABC$ è il baricentro del triangolo i cui vertici sono le proiezioni di $L$ sui lati $AB$, $BC$, $AC$. % lang-it
\end{proposition}

\begin{proposition}
[twierdzenie Schlömilcha] % lang-pl
[teorema di Schlömilch] % lang-it
    Trzy proste łączące środki boków trójkąta ze środkami odpowiednich wysokości są współpękowe, przechodzą przez punkt Lemoine'a. % lang-pl
    Le tre rette che congiungono i punti medi dei lati di un triangolo con i punti medi delle rispettive altezze sono concorrenti e passano per il punto di Lemoine. % lang-it
\end{proposition}

\index{twierdzenie!Schlömilcha} % lang-pl
\index{teorema!di Schlömilch} % lang-it

Napisze o tym duet Bogdańska, Neugebauer \cite[s. 195]{neugebauer_2018} (jako ćwiczenie, bez wzmianki o nazwie punktu), ale też Zetel \cite[s. 34]{zetel_2020}. % lang-pl
Ne scrivono Bogdańska e Neugebauer \cite[p. 195]{neugebauer_2018} (come esercizio, senza menzionare il nome del punto), ma anche Zetel \cite[p. 34]{zetel_2020}. % lang-it
% https://ems.press/content/serial-article-files/51635 VVV
Twierdzenie napotka w 1862 roku Oskar Schlömilch, cokolwiek to znaczy. % lang-pl
Il teorema fu scoperto da Oskar Schlömilch nel 1862, qualunque cosa ciò voglia dire. % lang-it

\begin{definition}
[antyrównoległa] % lang-pl
[antiparallela] % lang-it
    Antyrównoległa do boku $AB$ w trójkącie $ABC$ to prosta, która tnie bok naprzeciw wierzchołka $A$ pod takim samym kątem jak ten przy wierzchołku $A$ oraz bok naprzeciw wierzchołka $B$ pod takim samym kątem jak ten przy wierzchołku $B$. % lang-pl
    Un'antiparallela al lato $AB$ nel triangolo $ABC$ è una retta che interseca il lato opposto al vertice $A$ con lo stesso angolo presente in $A$ e il lato opposto al vertice $B$ con lo stesso angolo presente in $B$. % lang-it
\end{definition}

\begin{proposition}
    Trzy antyrównoległe do boków trójkąta przechodzące przez punkt Lemoine'a wyznaczają sześć nowych punktów na brzegu trójkąta. Leżą one na pierwszym okręgu Lemoine'a. % lang-pl
    Le tre antiparallele ai lati del triangolo passanti per il punto di Lemoine determinano sei nuovi punti sul contorno del triangolo. Essi giacciono sul primo cerchio di Lemoine. % lang-it
\end{proposition}

\begin{proposition}
    Trzy równoległe do boków trójkąta przechodzące przez punkt Lemoine'a wyznaczają sześć nowych punktów na brzegu trójkąta. Leżą one na drugim okręgu Lemoine'a. % lang-pl
    Le tre rette parallele ai lati del triangolo passanti per il punto di Lemoine determinano sei nuovi punti sul contorno del triangolo. Essi giacciono sul secondo cerchio di Lemoine. % lang-it
\end{proposition}

Istnieje jeszcze trzeci okrąg Lemoine'a, ale jego konstrukcja nie jest już taka prosta. % lang-pl
Esiste anche un terzo cerchio di Lemoine, ma la sua costruzione non è altrettanto semplice. % lang-it

%

\todofoot{Czewiany, w tym izotomiczne i izogonalne.}

\todofoot{Twierdzenie Steinera}

% 
%

\section{Twierdzenie van Aubela}
O twierdzeniu van Aubela pisze Zetel \cite[s. 24]{zetel_2020}.

% 

% \subsection{Twierdzenie Carnota???}
% \begin{proposition}[twierdzenie Carnota???]
% Neugebauer, strona 108.
% \end{proposition}

% biegun i biegunowa
%

\section{Twierdzenie Desargues'a} % lang-pl
\section{Teorema di Desargues} % lang-it

\begin{proposition}
[twierdzenie Desargues'a] % lang-pl
[teorema di Desargues] % lang-it
\index{twierdzenie!Desargues'a}% % lang-pl
\index{teorema!di Desargues}% % lang-it
    Niech $ABC$ i $A'B'C'$ będą dwoma trójkątami o parami różnych wierzchołkach, przy czym żaden wierzchołek jednego trójkąta nie leży na prostej zawierającej bok drugiego. % lang-pl
    Wtedy proste $AA'$, $BB'$, $CC'$ są współpękowe (mówimy, że trójkąty są \emph{perspektywiczne względem punktu}) wtedy i tylko wtedy, gdy punkty przecięcia odpowiednich boków $AB \cdot A'B'$, $BC \cdot B'C'$, $CA \cdot C'A'$ są współliniowe (mówimy, że trójkąty są \emph{perspektywiczne względem prostej}) -- patrz Neugebauer, strona 109. % lang-pl
    Sia $ABC$ e $A'B'C'$ due triangoli con vertici a due a due distinti, tali che nessun vertice dell'uno giaccia sulla retta contenente un lato dell'altro. % lang-it
    Allora le rette $AA'$, $BB'$, $CC'$ sono concorrenti (diciamo che i triangoli sono \emph{prospettivi rispetto a un punto}) se e solo se i punti di intersezione dei lati corrispondenti $AB \cdot A'B'$, $BC \cdot B'C'$, $CA \cdot C'A'$ sono allineati (diciamo che i triangoli sono \emph{prospettivi rispetto a una retta}) -- si veda Neugebauer, pagina 109. % lang-it
    % TODO: https://en.wikipedia.org/wiki/Desargues%27s_theorem

    Eves s. 245: Desargues two triangle theorem for the plane: copolar triangles are coaxial % lang-pl
    Eves p. 245: teorema dei due triangoli di Desargues nel piano: triangoli copolari sono coassiali % lang-it
\end{proposition}
% 1. Zna pojęcie inwolucji rzutowych.  Zna i potrafi stosować twierdzenia inwolucyjne Desarguesa.

Piszą o nim Audin \cite[s. 26, 151]{audin_2003}, Hartshorne w formie ćwiczenia \cite[s. 183]{hartshorne_2000}. % lang-pl
Eves \cite[s. 362]{eves1_1972} poda przykład Foresta Raya Moultona z 1902 roku płaszczyzny rzutowej, na której nie jest spełnione (!). % lang-pl
Ne parlano Audin \cite[p. 26, 151]{audin_2003} e Hartshorne sotto forma di esercizio \cite[p. 183]{hartshorne_2000}. % lang-it
Eves \cite[p. 362]{eves1_1972} fornirà un esempio di piano proiettivo del 1902, dovuto a Forest Ray Moulton, nel quale esso non è soddisfatto (!). % lang-it
\index[persons]{Moulton, Forest Ray}% https://en.wikipedia.org/wiki/Forest_Ray_Moulton
% Eves 365: Desaurgian/Pappian/classic

https://www.deltami.edu.pl/2015/04/desargues-i-nozyce/

https://www.deltami.edu.pl/2015/03/dowod-w-stylu-greckim/

%
% eksperymentalne tłumaczenie na włoski

%

\section{Twierdzenie Pappusa} % lang-pl
\section{Teorema di Pappo} % lang-it

\begin{proposition}
[twierdzenie Pappusa] % lang-pl
[teorema di Pappo] % lang-it
\index{twierdzenie!Pappusa}% % lang-pl
	Dane są dwie różne proste, $k$ oraz $l$ i sześć punktów różnych od $k \cdot l$: $K_1$, $K_2$, $K_3$ na prostej $k$, $L_1$, $L_2$, $L_3$ na prostej $l$. % lang-pl
	Wtedy punkty $K_1L_1 \cdot K_3L_2$, $L_1K_2 \cdot K_3L_3$ oraz $K_2L_2 \cdot K_1L_3$ są współliniowe. % lang-pl
\index{teorema!di Pappo}% % lang-it
	Siano date due rette distinte $k$ e $l$ e sei punti distinti da $k \cdot l$: $K_1$, $K_2$, $K_3$ sulla retta $k$, $L_1$, $L_2$, $L_3$ sulla retta $l$. % lang-it
	Allora i punti $K_1L_1 \cdot K_3L_2$, $L_1K_2 \cdot K_3L_3$ e $K_2L_2 \cdot K_1L_3$ sono allineati. % lang-it
\end{proposition}

% Eves s. 250 

% TODO: https://en.wikipedia.org/wiki/Pappus%27s_hexagon_theorem

Piszą o nim Audin \cite[s. 25, 151, 171]{audin_2003} (w wersji afinicznej, potem rzutowej), Neugebauer \cite[s. 115, 116]{neugebauer_2018}, Hartshorne \cite[s. 131, 132]{hartshorne2000} (wersja afinicza). % lang-pl
Ne parlano Audin \cite[pp. 25, 151, 171]{audin_2003} (prima nella versione affine, poi in quella proiettiva), Neugebauer \cite[pp. 115, 116]{neugebauer_2018} e Hartshorne \cite[pp. 131, 132]{hartshorne2000} (versione affine). % lang-it
% Hartshorne 62?
Hartshorne wspomni, że Hilbert pokazał w 1971, że niekoniecznie przemienne ciało\footnote{Czyli pierścień z dzieleniem.} $F$ jest przemienne wtedy i tylko wtedy, gdy twierdzenie Pappusa zachodzi w płaszczyźnie $F \times F$. % lang-pl
Hartshorne osserva che Hilbert mostrò nel 1971 che un corpo non necessariamente commutativo\footnote{Cioè un anello con divisione.} $F$ è commutativo se e solo se il teorema di Pappo vale nel piano $F \times F$. % lang-it

Warto zwrócić uwagę na podobieństwa między twierdzeniami Pascala i Pappusa. % lang-pl
Vale la pena osservare le somiglianze tra i teoremi di Pascal e di Pappo. % lang-it

% \item Zna przykłady przekształceń rzutowych i umie je stosować w zadaniach i dowodach twierdzeń rzutowych (Desarguesa, Pappusa, Pascala, Brianchona). zna pojęcia: biegun i biegunowa i potrafi formułować twierdzenia dualne.  
% wg Neugebauera, Brianchon jest dualny do Pascala

% Pascal, Brianchon, Chasles -> Eves s. 256 

https://www.deltami.edu.pl/2015/03/dowod-w-stylu-greckim/

%
% eksperymentalne tłumaczenie na włoski

%% src/chapters/07-wspolliniowosc/pascal_brianchon.tex
%

\section{Twierdzenie Pascala i Brianchona} % lang-pl

\section{Teoremi di Pascal e Brianchon} % lang-it

\begin{proposition}[twierdzenie Pascala] % lang-pl
	Neugebauer, strona 113. % lang-pl
\begin{proposition}[teorema di Pascal] % lang-it
	
	Neugebauer, pagina 113. % lang-it
	% TODO: https://en.wikipedia.org/wiki/Pascal%27s_theorem
	\index{twierdzenie!Pascala} % lang-pl
	
	\index{teorema!di Pascal} % lang-it
\end{proposition}

Audin \cite[s. 103, 107, 209]{audin_2003} przytoczy warianty tego twierdzenia. % lang-pl
Audin \cite[pp. 103, 107, 209]{audin_2003} % lang-it
presenta varie versioni di questo teorema. % lang-it

Joanna Jaszuńska poda je jako ćwiczenie w $\Delta_{16}^{11}$. % lang-pl
Joanna Jaszuńska lo propone come esercizio in $\Delta_{16}^{11}$. % lang-it
% https://www.deltami.edu.pl/2016/11/dwa-trojkaty/
% todo: i w 2014/9
Inna nazwa to \emph{hexagrammum mysticum theorem}. % lang-pl

Un altro nome è \emph{hexagrammum mysticum theorem}. % lang-it
% Pascal's theorem is the polar reciprocal and projective dual of Brianchon's theorem.
% It was formulated by Blaise Pascal in a note written in 1639 when he was 16 years old and published the following year as a broadside titled "Essay pour les coniques. Par B. P."[1]
% Pascal's theorem is a special case of the Cayley–Bacharach theorem.
% The converse is the Braikenridge–Maclaurin theorem, named for 18th-century British mathematicians William Braikenridge and Colin Maclaurin (Mills 1984), which states that if the three intersection points of the three pairs of lines through opposite sides of a hexagon lie on a line, then the six vertices of the hexagon lie on a conic; the conic may be degenerate, as in Pappus's theorem.[2] The Braikenridge–Maclaurin theorem may be applied in the Braikenridge–Maclaurin construction, which is a synthetic construction of the conic defined by five points, by varying the sixth point.

% As Thomas Kirkman proved in 1849, these 60 lines can be associated with 60 points in such a way that each point is on three lines and each line contains three points. The 60 points formed in this way are now known as the Kirkman points.[5] The Pascal lines also pass, three at a time, through 20 Steiner points. There are 20 Cayley lines which consist of a Steiner point and three Kirkman points. The Steiner points also lie, four at a time, on 15 Plücker lines. Furthermore, the 20 Cayley lines pass four at a time through 15 points known as the Salmon points.[6]

\begin{theorem}[Pascala] % lang-pl
	Jeżeli punkty $p_1$, $p_2$, $p_3$, $p_4$, $p_5$, $p_6$ leżą na pewnej stożkowej, to punkty $p_1p_2 \cdot p_4p_5$, $p_2p_3 \cdot p_5p_6$ oraz $p_3p_4 \cdot p_6p_1$ są współliniowe. % lang-pl
\begin{theorem}[di Pascal] % lang-it
	
	\index{twierdzenie!Pascala} % lang-pl
	Se i punti $p_1$, $p_2$, $p_3$, $p_4$, $p_5$, $p_6$ giacciono su una stessa conica, allora i punti $p_1p_2 \cdot p_4p_5$, $p_2p_3 \cdot p_5p_6$ e $p_3p_4 \cdot p_6p_1$ sono allineati. % lang-it
	
	\index{teorema!di Pascal} % lang-it
\end{theorem}

https://www.deltami.edu.pl/2014/09/twierdzenie-pascala/

\begin{theorem}[Brianchona] % lang-pl
	Jeżeli proste $l_1$, $l_2$, $l_3$, $l_4$, $l_5$, $l_6$ są styczne do pewnej stożkowej, to proste $p = (l_1 \cdot l_2)(l_4 \cdot l_5)$, $q = (l_2 \cdot l_3)(l_5 \cdot l_6)$ oraz $r = (l_3 \cdot l_4)(l_6 \cdot l_1)$ są współpękowe. % lang-pl
\begin{theorem}[di Brianchon] % lang-it
	
	\index{twierdzenie!Brianchona} % lang-pl
	Se le rette $l_1$, $l_2$, $l_3$, $l_4$, $l_5$, $l_6$ sono tangenti a una stessa conica, allora le rette $p = (l_1 \cdot l_2)(l_4 \cdot l_5)$, $q = (l_2 \cdot l_3)(l_5 \cdot l_6)$ e $r = (l_3 \cdot l_4)(l_6 \cdot l_1)$ sono concorrenti. % lang-it
	
	\index{teorema!di Brianchon} % lang-it
\end{theorem}

%  Coxeter, H. S. M. (1987). Projective Geometry (2nd ed.). Springer-Verlag. Theorem 9.15, p. 83. ISBN 0-387-96532-7.
% The polar reciprocal and projective dual of this theorem give Pascal's theorem.

To sformułowanie pojawi się u Bogdańskiej, Neugebauera \cite[s. 265, 266]{neugebauer_2018}. % lang-pl

Napisze o nim także nieznany autor w $\Delta_{84}^{11}$, Jaszuńska w $\Delta_{12}^4$, Eves \cite[s. 143]{eves1_1972}, a w prostszej wersji także Guzicki \cite[s. 238]{guzicki_2021}. % lang-pl
Questa formulazione compare in Bogdańska e Neugebauer \cite[pp. 265, 266]{neugebauer_2018}. % lang-it

Ne scrivono inoltre un autore anonimo in $\Delta_{84}^{11}$, Jaszuńska in $\Delta_{12}^4$, Eves \cite[p. 143]{eves1_1972} e, in una versione più semplice, anche Guzicki \cite[p. 238]{guzicki_2021}. % lang-it
% https://www.deltami.edu.pl/2012/04/twierdzenie-brianchona/

%
\todofoot{Twierdzenie Kirkmana: jeśli część wspólna dwóch trójkątów wpisanych w okrąg jest sześciokątem wypukłym, to główne przekątne tego sześciokąta przecinają się w jednym punkcie.} % lang-pl
\index{twierdzenie!Kirkmana} % lang-pl
\todofoot{Teorema di Kirkman: se l'intersezione di due triangoli inscritti in un cerchio è un esagono convesso, allora le diagonali principali di tale esagono si incontrano in un unico punto.} % lang-it

\index{teorema!di Kirkman} % lang-it

% \section{Różności}
\section{Twierdzenie Hirotaki?} % lang-pl

\section{Teorema di Hirotaka?} % lang-it
% eksperymentalne tłumaczenie na włoski

%

Twierdzenie Hirotaki: dany jest czworokąt cykliczny, wówczas proste są współpękowe. % lang-pl
\index{twierdzenie!Hirotaki} % lang-pl
Teorema di Hirotaka: dato un quadrilatero ciclico, allora le rette sono concorrenti. % lang-it
\index{teorema!Hirotaka} % lang-it
% Hirotaka pogrubiony u Neugebauera, strona 45

%

%